\documentclass[openany,fontset=fandol]{ctexbook}
\usepackage{amsmath,amssymb,amsthm}
\usepackage[unicode,hidelinks]{hyperref}
\setlength{\emergencystretch}{2em}
\raggedbottom
\DeclareMathOperator{\Spec}{Spec}
\newcommand{\PP}{\mathbf P}
\newcommand{\GG}{\mathcal G}
\newcommand{\dd}{\delta}
\newcommand{\tagitem}[2]{\par\medskip\noindent\textbf{#1\quad #2}\enspace}
\begin{document}
\title{代数几何基础第五卷预备笔记\\法文覆盖审计：截面的连通性与开性}
\author{第93批\quad 法、英、中文三方校订}
\date{2026年11月23日}
\maketitle

\section*{范围与接续}
逐段对读法文扫描物理19页下部至21页5.2前，英编稿物理15页的4.4至16页的5.1。
本批含推论4.4、作者补记及5.1，另记录法20页的废稿层。第三辑已有旧译，
本批是待替换校订稿，未改动1158页合订本。数学补证与原译文分开。

\section*{20.4（续）\quad “充分一般”的截面}
\tagitem{推论}{4.4}
考虑如下性质 $(C_m)$：$\overline Z$ 去掉任意维数不超过 $m$ 的闭子集后仍连通，
其中 $\overline Z=Z\times_K\overline K$，$\overline K$ 为 $K$ 的一个代数闭包。
令 $E$ 为使 $Y_\xi$ 在 $k(\xi)$ 上满足 $(C_m)$ 的 $\xi\in P^\vee$ 所组成的集合。
\begin{enumerate}
\item[(a)] $E$ 可构造。
\item[(b)] 设 $S=\Spec k$，$k$ 为域，并设 $f:X\to P$ 拟有限。
  若 $X$ 的每个不可约分支 $X_i$ 都满足 $\dim\overline{f(X_i)}\geq2$，
  且 $X$ 在 $k$ 上满足 $(C_{m+1})$，则 $E$ 包含 $P^\vee$ 的稠密开集 $U$。\footnote{“$f$ 拟有限”在本页陈述中未列，正文按下一页作者的更正加入。法第19页明确写 $f(X_i)$，英稿误为 $f(X)$；法稿末行以 $P$ 指参数空间，未辨得对偶标记，英稿同处亦作 $P$。由 $E$ 的定义及4.1，正文为 $P^\vee$。英稿末行的 $E_\xi$ 这里也还原成 $E$。}
\end{enumerate}

\noindent\textbf{作者补记。}
可构造性由 AQT 得到（第9段其实漏了这个事实，也许还来得及补上）。
(b)原则上来自3.3，不过3.3的表述过于特殊，须相应推广；所给证明基本不变。
另外，4.4 的表述有一处笔误：\emph{只有当 $f$ 拟有限时，原表述才适用}。
在一般情形，不应考察 $X$、$Y_\xi$ 中闭子集本身的维数，
而应考察它们在 $P$、$H_\xi$ 中的像的维数。%
\footnote{法第20页是限制式 \emph{ne\ldots que si}，即“仅当……”。英稿把它翻成了“若 $f$ 拟有限，则原式不适用”，方向相反。原稿未展开 AQT，英编者脚注也询问其含义；此处保留不确定性。}
\textit{Redactor demerdetur.}\footnote{保留原文末尾的戏拟拉丁语。英编稿同样留待翻译，本批不作确译。}

\section*{5\quad Seidenberg 型定理}
\tagitem{5.1}{}
在本节中，我们给出保证4.1所定义的 $E$ 为开集的条件。\footnote{法第21页在“给出条件”后写省略号，注明“见前页误划掉的文字”。故以下据第20页相应段复原，并与英编稿5.1互核。旧号4.5的这段虽被划销，却由5.1明确召回；不能将同页的旧命题4.6也按此处理。}
这里研究局部性质 $\mathcal P$，以便定义 $Y$ 中的子集 $U$：
若 $y\in Y$ 在 $P^\vee$ 中的像为 $\xi$，则 $y\in U$ 当且仅当
$Y_\xi$ 在 $y$ 处满足 $\mathcal P$，或在模的版本中 $\GG_\xi$ 在 $y$ 处满足 $\mathcal P$。
先用第12段给出 $U$ 为开的判据。\footnote{英编稿确认指 EGA IV，并补明 $E$ 的参数丛为 $P^\vee$。原文在这里只说明后面的论证路线，不是在任意情况下已经证明 $U$ 为开。}
由
\[
 E=P^\vee\setminus h(Y\setminus U),\qquad h:Y\longrightarrow P^\vee,
\]
可见，若 $U$ 是开集且 $X$ 在 $S$ 上固有，则 $E$ 也是开集：
因为 $Y$ 是 $X\times_S P^\vee$ 中的闭子概形，故 $h$ 固有，因而是闭映射。

\section*{废稿层}
法第20页下半另有被交叉划销的旧命题4.6，尚可辨认“有限表示”、
“相对于 $S$ Cohen--Macaulay”、“$d-1$”及“在 $P^\vee$ 上平坦”等片段。
原句有多重覆盖和截断，法第21页重起第5节，没有像5.1那样召回旧4.6。
不把它拼接成定稿或证明；原扫描仍附在工程中。

\section*{译校补证：为什么要量像的维数}
在一个代数闭底域上，对闭子集 $W\subset X$ 记
$\dd_f(W)=\dim\overline{f(W)}$，并置 $\dd_f(\varnothing)=-\infty$。
换底到代数闭包后能否被一个 $\dd_f\leq m$ 的闭子集断开，
是一般态射的\emph{像维数版本}。这里证明它在泛截面上的传递，
不以此替代4.4(a)的可构造性论证。

\paragraph{拟有限情形。}
若 $f$ 拟有限，对任一不可约闭子集 $W$，其泛点在像闭包泛点上方的剩余域扩张有限。
所以两个剩余域在底域上的超越次数相同，得到
$\dim W=\dim\overline{f(W)}$。对可约的 $W$ 取各分支维数的最大值。
拟有限性对基变换和限制到截面保持，故两种维数判据对 $X$、$Y_\xi$ 均相同。

\paragraph{图判据。}
设 $X$ 的不可约分支为 $X_i$。给 $\dd_f(X_i\cap X_j)>a$ 的两分支连一条边。
若图不连通，将顶点划分为两个非空、没有跨边的部分。删去全部跨部交集的并；
这是一闭集，$\dd_f\leq a$。两侧仍非空，因为一个不可约分支不包含在其他分支的有限并中。
两侧是互不相交的相对闭集，故删去后不连通。
反之，若图连通，删去任意 $\dd_f\leq a$ 的闭集，不会删掉任何一条边的全部交集；
每个留下的分支是不可约开子集，所以由边连接的并仍连通。
单分支时用不可约性直接判断，空集按拓扑学惯例视为连通。

\paragraph{泛截面。}
可以在 $\overline k$ 上做图判据，并用同时包含 $\overline k$ 和 $k(\eta)$
的 $k(\eta)$ 的代数闭包判定截面的几何性质。每个 $k$ 上的不可约分支上方，
$\overline k$ 上各分支由 Galois 作用传递置换（可分闭包到代数闭包的纯不可分扩张
不改变分支的拓扑），故它们的像维数相同，而且这个共同维数就是原分支的像维数。
因此假设的像维数下界也适用于每个几何分支。按3.1，每个几何分支的泛截面
在再扩到代数闭包后仍不可约。

对任意闭集 $W$，像闭包的最高维分支若维数为 $d\geq1$，
则泛超平面使这个像维数恰降为 $d-1$：
每个最高维分支中都有包含在 $f(W)$ 里的稠密开子集，因为像是可构造集
（参见 \href{https://stacks.math.columbia.edu/tag/054K}{Chevalley 定理}）。
按2.3、2.6，泛超平面与其闭包相交为纯 $d-1$ 维，
而补集的闭包维数至多为 $d-1$，其截面至多为 $d-2$；
零维部分泛截面不遇到。取闭包后，$W$ 的截面的像因此仍为 $d-1$ 维。
这些维数不会因进一步扩域而改变。关于几何分支的 Galois 作用，参见
\href{https://stacks.math.columbia.edu/tag/0364}{Stacks 项目第33.8节}。

现在若 $X_{\overline k}$ 不能被一个 $\dd_f\leq m+1$ 的闭集断开，
图判据给出一张连通图，每条边的交集像维数至少为 $m+2$。
截面后同一交集的像维数至少为 $m+1$。
因此删去 $\dd_{f_\eta}\leq m$ 的闭集仍留有这些连接，
各分支仍由原连通图串起来。这就证明了所需的泛截面断言。
若相应的参数集合已知可构造，再据4.1推出它含稠密开集。
这里没有将 AQT 认作某一已辨明的定理，也未独立证明任意底上的可构造性。

\paragraph{一般态射不能直接保留原维数：一个反例。}
取代数闭域上的 $X=V(uv)\subset\mathbf A^4_{u,v,w,z}$，
以及 $f:X\to\PP^2$，$f=[1:u+v:(u+v)w]$。
两个三维分支相交于平面 $u=v=0$；删去维数至多1的闭集后，
两分支的剩余部分仍不可约且交集非空，故 $X$ 满足 $(C_1)$。
每个分支的像都在仿射图的第一坐标非零处包含稠密开集，故像闭包为二维；
但 $f$ 把交平面压为一点。泛截面方程
\[
 a+(u+v)(b+cw)=0
\]
中 $a\ne0$；它避开交平面，两分支各自的截面都非空
（例如先选 $w$ 使 $b+cw\ne0$）。所以泛截面分成两个不相交的非空闭部，
连不删点都不连通，因而不满足 $(C_0)$。
这不是拟有限版本的反例；它说明为什么一般情形须改用像维数。

\medskip
\noindent\textbf{接续点。}
法物理21页5.2、英物理16页5.2，先核可逆模的局部平凡化（不是任意 $\mathcal F$ 与 $\mathcal O_X$ 局部同构）
及“除子”与有效 Cartier 除子的区分。法22的5.3手改将有限表示及 $S$-平坦放在(a)，
(b)的 $s\in P^\vee$ 疑应为 $s\in S$。这些还未校入正文。
正式 EGA IV.1 仍另从物理221页提要及223页定义（1.1.1）接续。
\end{document}
